\chapter{बहु-इलेक्ट्रॉन परमाणु और पद प्रतीक}\label{ch:b3:many-electron-atoms}

1868 में सूर्य के स्पेक्ट्रम की एक पीली रेखा ने एक ऐसे तत्व को प्रकट किया जो
पृथ्वी पर अज्ञात था, हीलियम। जब अंततः 1895 में उसे प्रयोगशाला में पृथक किया गया,
तो उसका स्पेक्ट्रम \emph{दो} तत्वों के स्पेक्ट्रम जैसा दिखा: रेखाओं के दो कुल,
जिन्हें कुछ समय तक ``ऑर्थोहीलियम'' और ``पैराहीलियम'' का माना गया, और जिन्होंने
आपस में कभी प्रकाश का विनिमय नहीं किया। हीलियम एक ही है। दोनों कुल
एक ही परमाणु की त्रिक और \omterm{def:b3:many-electron-atoms:singlet-triplet}{एकल अवस्थाएँ} हैं, जिन्हें इलेक्ट्रॉन का चक्रण और किसी भी
बल से अधिक मौलिक एक नियम अलग रखते हैं: दो इलेक्ट्रॉनों का तरंग फलन उनके परस्पर
विनिमय पर चिह्न बदलता है। यह अध्याय उस नियम का अनुसरण पाउली सिद्धांत से लेकर
उन \omterm{def:b3:many-electron-atoms:term}{पद प्रतीकों} तक करता है जिनसे स्पेक्ट्रमिकीविद् हर परमाणु और आयन की हर अवस्था
को अंकित करते हैं --- उन संक्रमण-धातु आयनों सहित जिनके रंगों की व्याख्या
\cref{ch:b3:complex-spectra-magnetism} में की गई है।

\begin{recall}
प्रथम वर्ष के खंड ने इलेक्ट्रॉनों को चार क्वांटम संख्याओं $n$, $l$,
$m_l$, $m_s$ से अंकित किया, पाउली सिद्धांत, क्लेचकोव्स्की नियम और हुंड के नियम
से मूल विन्यास बनाए, और अयुग्मित इलेक्ट्रॉन गिने। द्वितीय वर्ष के
खंड ने कक्षक ऊर्जाएँ और आवरण के लिए स्लेटर के नियम दिए।
\Cref{ch:b3:quantum-model-systems} ने \omterm{def:b3:quantum-model-systems:operator}{संकारक}, \omterm{def:b3:quantum-model-systems:operator}{अभिलक्षणिक मान},
घूर्णक का कोणीय संवेग ($\hat L^2$, \omterm{def:b3:quantum-model-systems:operator}{अभिलक्षणिक मानों} $l(l+1)\hbar^2$ के साथ) और
हाइड्रोजन परमाणु दिए।
\end{recall}

\section{चक्रण और चक्रण-कक्षक}

इलेक्ट्रॉन एक नैज कोणीय संवेग, अपना चक्रण, वहन करता है, जिसकी क्वांटम
संख्या $s = \frac12$ है: $\hat S^2$ का एकमात्र \omterm{def:b3:quantum-model-systems:operator}{अभिलक्षणिक मान} $s(s+1)\hbar^2 =
\frac34\hbar^2$ है, और $\hat S_z$ के दो \omterm{def:b3:quantum-model-systems:operator}{अभिलक्षणिक मान} $m_s\hbar = \pm\frac12\hbar$।
दो चक्रण फलन $\alpha$ ($m_s = +\frac12$) और $\beta$
($m_s = -\frac12$) लिखे जाते हैं; वे प्रसामान्य-लांबिक हैं, $\langle\alpha|\alpha\rangle =
\langle\beta|\beta\rangle = 1$ और $\langle\alpha|\beta\rangle = 0$। चक्रण का कोई
चिरसम्मत प्रतिरूप नहीं है; वह किसी छोटे गोले का घूर्णन नहीं है, और
अनापेक्षिकीय \omterm{def:b3:quantum-model-systems:schrodinger}{हैमिल्टनी संकारक} में वह बिल्कुल प्रवेश नहीं करता।

\begin{definition}[चक्रण-कक्षक]\label{def:b3:many-electron-atoms:spin-orbital}
\emph{चक्रण-कक्षक}\index{चक्रण-कक्षक} किसी आकाशी कक्षक
$\phi(\mathbf r)$ और एक चक्रण फलन का गुणनफल है: $\phi\alpha$ या $\phi\beta$। एक
चक्रण-कक्षक एक इलेक्ट्रॉन का पूरा वर्णन धारण करता है।
\end{definition}

इस प्रकार प्रथम वर्ष के खंड में $(n, l, m_l)$ से अंकित परमाणु कक्षक दो
\omterm{def:b3:many-electron-atoms:spin-orbital}{चक्रण-कक्षक} देता है, इसीलिए एक उपकोश में $2(2l+1)$ इलेक्ट्रॉन आते हैं। वास्तविक
प्रश्न यह है कि एक \omterm{def:b3:many-electron-atoms:spin-orbital}{चक्रण-कक्षक} में अधिकतम एक ही इलेक्ट्रॉन क्यों आता है।

\section{अविभेद्य इलेक्ट्रॉन और स्लेटर सारणिक}

\begin{definition}[अविभेद्य कण, प्रतिसममित तरंग फलन]\label{def:b3:many-electron-atoms:indistinguishable}
कण \emph{अविभेद्य}\index{अविभेद्य कण} हैं यदि
कोई भी मापन उन्हें अलग न पहचान सके: हर इलेक्ट्रॉन हर दूसरे के सर्वसम है।
अनेक इलेक्ट्रॉनों का तरंग फलन \emph{प्रतिसममित
तरंग फलन}\index{प्रतिसममित तरंग फलन} है यदि किन्हीं दो इलेक्ट्रॉनों के
निर्देशांकों (आकाश और चक्रण) का विनिमय करने पर वह चिह्न बदलता है:
$\Psi(\dots,i,\dots,j,\dots) = -\Psi(\dots,j,\dots,i,\dots)$।
\end{definition}

केवल अविभेद्यता यह माँगती है कि विनिमय से $|\Psi|^2$ अपरिवर्तित रहे,
अतः $\Psi$ अधिक से अधिक चिह्न बदल सकता है। प्रकृति ने चुनाव कर लिया है:
इलेक्ट्रॉनों के किसी भी निकाय का --- किन्हीं भी फ़र्मिऑनों का --- तरंग फलन
प्रतिसममित होता है। यह क्वांटम यांत्रिकी का एक अभिगृहीत है, जिसकी पुष्टि हर
परमाण्वीय स्पेक्ट्रम करता है।

\begin{definition}[स्लेटर सारणिक]\label{def:b3:many-electron-atoms:slater}
$N$ इलेक्ट्रॉनों के लिए, जो $N$ भिन्न \omterm{def:b3:many-electron-atoms:spin-orbital}{चक्रण-कक्षकों} $\chi_1, \dots, \chi_N$ में हैं,
\emph{स्लेटर सारणिक}\index{स्लेटर सारणिक} यह है:
\[
\Psi = \frac{1}{\sqrt{N!}}\begin{vmatrix}
\chi_1(1) & \chi_2(1) & \cdots & \chi_N(1)\\
\chi_1(2) & \chi_2(2) & \cdots & \chi_N(2)\\
\vdots & & & \vdots\\
\chi_1(N) & \chi_2(N) & \cdots & \chi_N(N)
\end{vmatrix},
\]
जहाँ पंक्ति $i$ में इलेक्ट्रॉन $i$ हर \omterm{def:b3:many-electron-atoms:spin-orbital}{चक्रण-कक्षक} में है। संक्षेप में इसे
$|\chi_1\chi_2\dots\chi_N|$ लिखा जाता है।
\end{definition}

\begin{theorem}[पाउली सिद्धांत]\label{thm:b3:many-electron-atoms:pauli}
\omterm{def:b3:many-electron-atoms:slater}{स्लेटर सारणिक} प्रतिसममित होता है, और यदि दो इलेक्ट्रॉन एक ही \omterm{def:b3:many-electron-atoms:spin-orbital}{चक्रण-कक्षक} में
हों तो वह लुप्त हो जाता है। अतः किसी परमाणु के कोई भी दो इलेक्ट्रॉनों की
चारों क्वांटम संख्याएँ समान नहीं हो सकतीं।
\end{theorem}

\begin{proof}
इलेक्ट्रॉनों $i$ और $j$ का विनिमय सारणिक की दो पंक्तियों का विनिमय करता है,
जिससे उसका चिह्न बदल जाता है: हर युग्म के लिए प्रतिसममिति लागू है। यदि दो
\omterm{def:b3:many-electron-atoms:spin-orbital}{चक्रण-कक्षक} बराबर हों, तो दो स्तंभ बराबर होते हैं और सारणिक शून्य है:
ऐसी कोई अवस्था नहीं होती। एक ही परमाणु के दो इलेक्ट्रॉन, जिनके $n$, $l$,
$m_l$ और $m_s$ समान हों, एक ही \omterm{def:b3:many-electron-atoms:spin-orbital}{चक्रण-कक्षक} में होते।
\end{proof}

\begin{example}[हीलियम की मूल अवस्था]\label{ex:b3:many-electron-atoms:helium-ground}
दोनों इलेक्ट्रॉन $1s$ में होने पर,
\[
\Psi = \frac{1}{\sqrt2}\begin{vmatrix}1s\alpha(1) & 1s\beta(1)\\ 1s\alpha(2) &
1s\beta(2)\end{vmatrix} = 1s(1)\,1s(2)\cdot\frac{1}{\sqrt2}\big[\alpha(1)\beta(2)
- \beta(1)\alpha(2)\big].
\]
आकाशी भाग सममित है, चक्रण भाग प्रतिसममित: दोनों चक्रण
युग्मित हैं, कुल चक्रण शून्य के साथ।
\end{example}

\section{हीलियम: कूलॉम प्रतिकर्षण और विनिमय}

हीलियम के एक इलेक्ट्रॉन को $2s$ में उत्तेजित कीजिए। हर कक्षक में एक इलेक्ट्रॉन के साथ
$1s\alpha$, $1s\beta$, $2s\alpha$, $2s\beta$ से चार सारणिक बनाए जा सकते हैं।
पुनर्व्यवस्थित करने पर वे एक आकाशी और एक चक्रण फलन के गुणनफल बन जाते हैं:
\[
\Psi_\pm = \frac{1}{\sqrt2}\big[1s(1)2s(2) \pm 2s(1)1s(2)\big]\times(\text{चक्रण
फलन}).
\]
प्रतिसममित आकाशी फलन $\Psi_-$ को किसी सममित
चक्रण फलन के साथ युग्मित होना होगा, और ऐसे तीन हैं: $\alpha(1)\alpha(2)$,
$\beta(1)\beta(2)$ और $\frac{1}{\sqrt2}[\alpha(1)\beta(2) + \beta(1)\alpha(2)]$, जो
तीन घटक $M_S = 1, 0, -1$ हैं, कुल चक्रण $S = 1$ के। सममित
$\Psi_+$ को एकमात्र प्रतिसममित चक्रण फलन के साथ युग्मित होना होगा, $S = 0$।

\begin{definition}[एकल और त्रिक अवस्थाएँ]\label{def:b3:many-electron-atoms:singlet-triplet}
कुल चक्रण $S = 0$ वाली अवस्था \emph{एकल अवस्था}\index{एकल अवस्था} है;
कुल चक्रण $S = 1$ वाली अवस्था, जिसके तीन घटक $M_S = 1, 0, -1$ चुंबकीय क्षेत्र की
अनुपस्थिति में समान ऊर्जा रखते हैं, \emph{त्रिक
अवस्था}\index{त्रिक अवस्था} है।
\end{definition}

\begin{definition}[विनिमय समाकल]\label{def:b3:many-electron-atoms:exchange}
दो कक्षकों $a$ और $b$ तथा इलेक्ट्रॉन प्रतिकर्षण $e^2/4\pi\varepsilon_0
r_{12}$ के लिए \emph{विनिमय समाकल}\index{विनिमय समाकल} यह है:
\[
K_{ab} = \iint a(1)b(2)\,\frac{e^2}{4\pi\varepsilon_0r_{12}}\,b(1)a(2)\,\dd\tau_1\dd\tau_2,
\]
यह वही समाकल है जो चिरसम्मत इलेक्ट्रॉन प्रतिकर्षण $J_{ab}$ (जिसमें
दोनों ओर $a(1)b(2)$ है), सिवाय इसके कि एक ओर इलेक्ट्रॉनों की अदला-बदली की गई है।
$K_{ab}$ धनात्मक है और उसका कोई चिरसम्मत प्रतिरूप नहीं है।
\end{definition}

\begin{proposition}[एकल और त्रिक ऊर्जाएँ]\label{prop:b3:many-electron-atoms:helium}
इलेक्ट्रॉन प्रतिकर्षण में प्रथम कोटि तक, हीलियम की $1s2s$ अवस्थाओं की
ऊर्जाएँ $E_\pm = E_{1s} + E_{2s} + J \pm K$ हैं, एकल ($+$)
त्रिक ($-$) से $2K$ ऊपर।
\end{proposition}

\begin{proof}
दोनों $\Psi_\pm$ की अक्षुब्ध ऊर्जा $E_{1s} + E_{2s}$ है (नाभिकीय आवेश 2 के
हाइड्रोजन-सदृश कक्षक)। प्रथम कोटि का संशोधन प्रसामान्यीकृत अवस्था में प्रतिकर्षण
$\hat V_{12}$ का \omterm{def:b3:quantum-model-systems:hermitian}{प्रत्याशा मान} है:
$\frac12\langle 1s(1)2s(2) \pm 2s(1)1s(2)|\hat V_{12}|1s(1)2s(2) \pm
2s(1)1s(2)\rangle$। दो प्रत्यक्ष पद $\frac12(J + J)$ देते हैं; दो तिर्यक
पद $\pm\frac12(K + K)$ देते हैं, क्योंकि $\hat V_{12}$ इलेक्ट्रॉनों में
सममित है। अतः $\Delta E = J \pm K$, और चक्रण फलन, जो प्रसामान्यीकृत हैं और
$\hat V_{12}$ से अप्रभावित, केवल 1 से गुणा करते हैं।
\end{proof}

त्रिक नीचे स्थित है यद्यपि \omterm{def:b3:quantum-model-systems:schrodinger}{हैमिल्टनी संकारक} में कोई चक्रण नहीं है: उसका
आकाशी फलन $\mathbf r_1 = \mathbf r_2$ होने पर लुप्त होता है, अतः दोनों
इलेक्ट्रॉन एक-दूसरे से बचते हैं और कम प्रतिकर्षण करते हैं। यह ``फ़र्मी छिद्र''
हुंड के पहले नियम का भौतिक सार है।

\begin{example}[हीलियम का विनिमय समाकल, मापित]\label{ex:b3:many-electron-atoms:K}
% ledger: lev:He.2s3S1, lev:He.2s1S0
हीलियम के $1s2s$ स्तर मूल अवस्था से \qty{159856}{cm^{-1}} (${}^3$S) और
\qty{166277}{cm^{-1}} (${}^1$S) ऊपर हैं: एकल
\qty{6421}{cm^{-1}}, अर्थात् \qty{0.796}{eV}, ऊँचा है, और $K(1s,2s) = \qty{0.398}{eV}$।
नीचे का चित्र दोनों कुल दिखाता है।
\end{example}

\section{रसेल--सॉन्डर्स युग्मन और पद प्रतीक}

किसी हल्के परमाणु में इलेक्ट्रॉनों के कक्षीय कोणीय संवेग जुड़कर एक
कुल $\mathbf L$ बनाते हैं, उनके चक्रण एक कुल $\mathbf S$, और केवल इसके बाद
$\mathbf L$ और $\mathbf S$ आपस में दुर्बल रूप से युग्मित होते हैं।

\begin{definition}[रसेल--सॉन्डर्स युग्मन]\label{def:b3:many-electron-atoms:russell-saunders}
\emph{रसेल--सॉन्डर्स युग्मन}\index{रसेल--सॉन्डर्स युग्मन} में
किसी विन्यास की अवस्थाएँ इनसे अंकित होती हैं: \emph{कुल कक्षीय कोणीय
संवेग क्वांटम संख्या}\index{कुल कक्षीय कोणीय संवेग क्वांटम संख्या}
$L$ ($\hat{\mathbf L}^2 = L(L+1)\hbar^2$, जहाँ $M_L = \sum m_l$), \emph{कुल
चक्रण क्वांटम संख्या}\index{कुल चक्रण क्वांटम संख्या} $S$ ($M_S = \sum m_s$),
और \emph{कुल कोणीय संवेग क्वांटम संख्या}\index{कुल कोणीय संवेग क्वांटम संख्या}
$J$, जो मान $|L - S|, \dots, L + S$ लेती है।
\end{definition}

\begin{definition}[स्पेक्ट्रमी पद, चक्रण बहुकता, पद प्रतीक]\label{def:b3:many-electron-atoms:term}
\emph{स्पेक्ट्रमी पद}\index{स्पेक्ट्रमी पद} किसी विन्यास की उन
$(2L+1)(2S+1)$ अवस्थाओं का समुच्चय है जिनके $L$ और $S$ दिए हुए हैं। उसकी
\emph{चक्रण बहुकता}\index{चक्रण बहुकता} $2S + 1$ है। \emph{पद प्रतीक}
\index{पद प्रतीक} ${}^{2S+1}L$ है, जहाँ अक्षर S, P, D, F, G, H
क्रमशः $L = 0, 1, 2, 3, 4, 5$ के लिए हैं; दिए हुए $J$ वाली अवस्था ${}^{2S+1}L_J$ लिखी जाती है ---
उदाहरण के लिए \termsym{3}{P}{2}।
\end{definition}

\begin{proposition}[पूर्ण कोश]\label{prop:b3:many-electron-atoms:closed-shell}
पूर्ण भरे उपकोश के लिए $L = 0$ और $S = 0$; वह \omterm{def:b3:many-electron-atoms:term}{पद प्रतीक} में कोई
योगदान नहीं देता।
\end{proposition}

\begin{proof}
पूर्ण भरे उपकोश में हर $m_l$ दो बार आता है और हर $m_s = \pm\frac12$
$2l+1$ बार, अतः $M_L = 0$ और $M_S = 0$; और उसे भरने का केवल एक ही
तरीका है (एक सारणिक)। $M_L = M_S = 0$ वाली एकमात्र अवस्था केवल
$L = 0$, $S = 0$ की हो सकती है।
\end{proof}

\begin{proposition}[किसी विन्यास की अवस्थाओं की गिनती]\label{prop:b3:many-electron-atoms:count}
$N$ इलेक्ट्रॉन, $2(2l+1)$ \omterm{def:b3:many-electron-atoms:spin-orbital}{चक्रण-कक्षकों} वाले उपकोश में,
$\binom{2(2l+1)}{N}$ सारणिक देते हैं, और उसके पदों की \omterm{def:b3:quantum-model-systems:degenerate}{अपभ्रष्टताओं} $(2L+1)(2S+1)$ का
योग यही संख्या होता है।
\end{proposition}

\begin{proof}
सारणिक अधिकृत \omterm{def:b3:many-electron-atoms:spin-orbital}{चक्रण-कक्षकों} के समुच्चय से निर्धारित होता है (उनका क्रम
केवल उसका चिह्न बदलता है): $N$ \omterm{def:b3:many-electron-atoms:spin-orbital}{चक्रण-कक्षक} चुने जाते हैं, कुल $2(2l+1)$ में से। पद वही
अवस्थाएँ हैं, पुनः समूहित, अतः उनकी गिनती मेल खानी चाहिए।
\end{proof}

\begin{method}[किसी विन्यास के पद]\label{met:b3:many-electron-atoms:terms}
\begin{enumerate}
  \item पाउली सिद्धांत द्वारा अनुमत सारणिकों को सूचीबद्ध कीजिए और उन्हें
        $M_L$ और $M_S$ के अनुसार सारणीबद्ध कीजिए।
  \item सबसे बड़ा $M_L$ ढूँढ़िए; उसके साथ आने वाले सबसे बड़े $M_S$ सहित
        वह $L = M_L$, $S = M_S$ वाला एक पद आरंभ करता है।
  \item हर उस युग्म $(M_L, M_S)$ के लिए एक सारणिक काट दीजिए जिसके लिए
        $|M_L| \le L$, $|M_S| \le S$।
  \item जो बचा है उसके साथ यही दोहराइए, जब तक सारणी खाली न हो जाए; गिनती जाँचिए।
\end{enumerate}
\end{method}

\begin{example}[$p^2$ के पद]\label{ex:b3:many-electron-atoms:p2}
$p$ ($m_l = 1, 0, -1$) में दो इलेक्ट्रॉन $\binom62 = 15$ सारणिक देते हैं।
$M_L$ और $M_S$ के अनुसार उनकी सारणी यह है:
\begin{center}\small
\begin{tabular}{c|ccc}
\hline
 & $M_S = 1$ & $M_S = 0$ & $M_S = -1$\\ \hline
$M_L = 2$ & -- & $(1^+,1^-)$ & --\\
$M_L = 1$ & $(1^+,0^+)$ & $(1^+,0^-)$, $(1^-,0^+)$ & $(1^-,0^-)$\\
$M_L = 0$ & $(1^+,-1^+)$ & $(1^+,-1^-)$, $(1^-,-1^+)$, $(0^+,0^-)$ & $(1^-,-1^-)$\\
$M_L = -1$ & $(0^+,-1^+)$ & $(0^+,-1^-)$, $(0^-,-1^+)$ & $(0^-,-1^-)$\\
$M_L = -2$ & -- & $(-1^+,-1^-)$ & --\\ \hline
\end{tabular}
\end{center}
($m_l^\pm$ से तात्पर्य $m_s = \pm\frac12$)। $M_L = 2$ केवल $M_S = 0$ के साथ आता है:
एक ${}^1$D पद (5 अवस्थाएँ)। शेष में सबसे बड़ा $M_L = 1$,
$M_S = 1$ के साथ आता है: एक ${}^3$P पद (9 अवस्थाएँ)। $M_L = M_S = 0$ वाली एक अवस्था बचती है: एक
${}^1$S पद। वास्तव में $5 + 9 + 1 = 15$।
\end{example}

\begin{example}[कार्बन, मापित]\label{ex:b3:many-electron-atoms:carbon}
% ledger: lev:C.3P1, lev:C.3P2, lev:C.1D2, lev:C.1S0
कार्बन का मूल विन्यास $2p^2$ स्तर \termsym{3}{P}{0},
\termsym{3}{P}{1}, \termsym{3}{P}{2} देता है, जो 0, 16.4 और \qty{43.4}{cm^{-1}} पर हैं, फिर
\termsym{1}{D}{2}, \qty{10193}{cm^{-1}} पर, और \termsym{1}{S}{0},
\qty{21648}{cm^{-1}} पर: ${}^3$P पद सबसे नीचे है, जैसा हुंड के नियम बताते हैं।
\end{example}

\begin{omfigure}
\begin{tikzpicture}
\begin{axis}[name=main, width=0.62\linewidth, height=6.4cm, xmin=-0.3, xmax=6.6,
  ymin=-1500, ymax=23500, axis y line=left, axis x line=none,
  ylabel={ऊर्जा / cm$^{-1}$}, unbounded coords=jump, clip=false,
  scaled y ticks=false, ytick={0,10000,20000},
  tick label style={font=\small}, label style={font=\small}]
\addplot[omstate] table[x=x, y=E] {figdata/out/bachelor-3-atomic-levels-carbon.dat};
\draw[densely dotted, gray] (axis cs:1,4858.5) -- (axis cs:2,29.6);
\draw[densely dotted, gray] (axis cs:1,4858.5) -- (axis cs:2,10192.7);
\draw[densely dotted, gray] (axis cs:1,4858.5) -- (axis cs:2,21648);
\draw[densely dotted, gray] (axis cs:3,10192.7) -- (axis cs:4,10192.7);
\draw[densely dotted, gray] (axis cs:3,21648) -- (axis cs:4,21648);
\node[font=\small, anchor=south] at (axis cs:0.5,4858.5) {$2p^2$};
\node[font=\small, anchor=south] at (axis cs:2.5,21648) {${}^1$S};
\node[font=\small, anchor=south] at (axis cs:2.5,10192.7) {${}^1$D};
\node[font=\small, anchor=south] at (axis cs:2.5,29.6) {${}^3$P};
\node[font=\small, anchor=west] at (axis cs:5.05,21648) {\termsym{1}{S}{0}};
\node[font=\small, anchor=west] at (axis cs:5.05,10192.7) {\termsym{1}{D}{2}};
\node[font=\scriptsize, anchor=north] at (axis cs:0.5,-600) {विन्यास};
\node[font=\scriptsize, anchor=north] at (axis cs:2.5,-600) {पद};
\node[font=\scriptsize, anchor=north] at (axis cs:4.5,-600) {स्तर};
\draw[gray, ->] (axis cs:3.05,60) -- (axis cs:6.9,2500);
\end{axis}
\begin{axis}[at={(main.south east)}, anchor=south west, xshift=1.2cm, yshift=0.9cm,
  width=3.6cm, height=3.6cm, xmin=-0.1, xmax=1.1, ymin=-5, ymax=50,
  axis y line=right, axis x line=none, unbounded coords=jump, ytick={0,16.4,43.4},
  yticklabels={0,16.4,43.4}, tick label style={font=\scriptsize}, clip=false,
  title={\scriptsize ${}^3$P, आवर्धित}, title style={yshift=-4pt}]
\addplot[omstate] table[x=x, y=E] {figdata/out/bachelor-3-atomic-levels-carbon-zoom.dat};
\node[font=\scriptsize, anchor=east] at (axis cs:0,0) {$J = 0$};
\node[font=\scriptsize, anchor=east] at (axis cs:0,16.4) {$J = 1$};
\node[font=\scriptsize, anchor=east] at (axis cs:0,43.4) {$J = 2$};
\end{axis}
\end{tikzpicture}

{\small कार्बन का मूल विन्यास: एक विन्यास, तीन पद
(उनकी मापी गई ऊर्जाओं पर; ${}^3$P पद अपने स्तरों के भारित माध्य पर,
विन्यास सभी पंद्रह अवस्थाओं के माध्य पर) और पाँच
स्तर। ${}^3$P का चक्रण--कक्षक विपाटन (आवर्धित,
\unit{cm^{-1}} में) पदों के बीच के अंतरालों से कुछ सौ गुना छोटा
है।}
\end{omfigure}

\section{हुंड के नियम और चक्रण--कक्षक युग्मन}

\begin{proposition}[पदों के लिए हुंड के नियम]\label{prop:b3:many-electron-atoms:hund-terms}
किसी परमाणु या आयन के मूल विन्यास के लिए, निम्नतम पद वह है जिसका
(1) $S$ सबसे बड़ा हो, फिर (2) उस $S$ के लिए $L$ सबसे बड़ा हो। उसके निम्नतम
स्तर का (3) $J = |L - S|$ होता है यदि उपकोश आधे से कम भरा हो, और $J = L +
S$ यदि वह आधे से अधिक भरा हो।
\end{proposition}
\begin{proof}[स्थिति]
ये नियम प्रायोगिक रूप से स्थापित हैं, केवल मूल विन्यासों के लिए;
ये प्रमेय नहीं हैं। नियम 1 पिछले अनुभाग का विनिमय स्थायीकरण है;
नियम 2 कहता है कि एक ही दिशा में घूमते इलेक्ट्रॉन कम बार मिलते हैं;
नियम 3 \omterm{def:b3:many-electron-atoms:spin-orbit}{चक्रण--कक्षक युग्मन} के चिह्न से निकलता है, जो आगे दिया गया है।
\end{proof}

\begin{method}[खाना आरेख से मूल पद]\label{met:b3:many-electron-atoms:ground-term}
\begin{enumerate}
  \item खुले उपकोश के खानों को $m_l = +l$ से नीचे की ओर भरिए, हर खाने में एक
        इलेक्ट्रॉन, समांतर चक्रणों के साथ, फिर $+l$ से पुनः युग्मित कीजिए।
  \item $S = \frac12 \times$(अयुग्मित इलेक्ट्रॉनों की संख्या); $L = |\sum m_l|$।
  \item आधे से कम भरे के लिए $J = |L - S|$, आधे से अधिक भरे के लिए $L + S$;
        ठीक आधे भरे के लिए $L = 0$ और $J = S$।
\end{enumerate}
\end{method}

\begin{example}[मूल पद]\label{ex:b3:many-electron-atoms:ground-terms}
% ledger: lev:Ti2.3P0, lev:Ni2.3P0, lev:Cr3.4P1_2
\ce{Ti^2+}, $d^2$: खाने $m_l = 2, 1$ भरे, $S = 1$, $L = 3$, $J = 2$:
\termsym{3}{F}{2}। \ce{Cr^3+}, $d^3$: $S = \frac32$, $L = 2 + 1 + 0 = 3$,
\termsym{4}{F}{3/2}। \ce{Fe^2+}, $d^6$: पाँच ऊपर, एक नीचे $m_l = 2$ में:
$S = 2$, $L = 2$, आधे से अधिक भरा, \termsym{5}{D}{4}। \ce{Ni^2+}, $d^8$:
$S = 1$, $L = 3$, \termsym{3}{F}{4}। हर एक उस मूल स्तर से मेल खाता है जिसे
परमाण्वीय स्पेक्ट्रमिकी मापती है।
\end{example}

\begin{definition}[चक्रण--कक्षक युग्मन, सूक्ष्म संरचना]\label{def:b3:many-electron-atoms:spin-orbit}
\emph{चक्रण--कक्षक युग्मन}\index{चक्रण--कक्षक युग्मन} किसी इलेक्ट्रॉन के
चक्रण और उसकी कक्षीय गति के चुंबकीय आघूर्णों के बीच की अन्योन्यक्रिया है, जिसे किसी
पद के लिए $A\,\hat{\mathbf L}\cdot\hat{\mathbf S}/\hbar^2$ लिखा जाता है। यह
किसी पद को भिन्न $J$ वाले उसके स्तरों में विपाटित करता है: यही पद की \emph{सूक्ष्म
संरचना}\index{सूक्ष्म संरचना} है।
\end{definition}

\begin{proposition}[लांडे का अंतराल नियम]\label{prop:b3:many-electron-atoms:lande}
किसी पद के भीतर $E(J) = E_0 + \frac{A}{2}[J(J+1) - L(L+1) - S(S+1)]$, अतः
$E(J) - E(J - 1) = AJ$। आधे से कम भरे उपकोश के लिए $A > 0$ (सामान्य),
आधे से अधिक भरे के लिए $A < 0$ (प्रतिलोमित)।
\end{proposition}

\begin{proof}
$\hat{\mathbf J} = \hat{\mathbf L} + \hat{\mathbf S}$ से $\hat{\mathbf J}^2 =
\hat{\mathbf L}^2 + \hat{\mathbf S}^2 + 2\hat{\mathbf L}\cdot\hat{\mathbf S}$, अतः
$\hat{\mathbf L}\cdot\hat{\mathbf S} = \frac12(\hat{\mathbf J}^2 - \hat{\mathbf L}^2 -
\hat{\mathbf S}^2)$, जिसका दिए हुए $J$, $L$, $S$ वाली अवस्था में \omterm{def:b3:quantum-model-systems:operator}{अभिलक्षणिक मान}
$\frac{\hbar^2}{2}[J(J+1) - L(L+1) - S(S+1)]$ है। $J$ और
$J - 1$ के बीच का अंतर $\frac A2[J(J+1) - (J-1)J] = AJ$ है। $A$ का चिह्न स्वीकार किया गया है:
आधे से अधिक भरा उपकोश संगत संख्या के धनात्मक
छिद्रों की तरह व्यवहार करता है।
\end{proof}

\begin{example}[अंतराल नियम का परीक्षण]\label{ex:b3:many-electron-atoms:lande-test}
% ledger: lev:C.3P1, lev:C.3P2, lev:O.3P1, lev:O.3P0
${}^3$P के लिए नियम अनुपात $J = 2 : 1$, अर्थात् 2, में अंतराल बताता है।
कार्बन: $(43.4 - 16.4)/16.4 = 1.65$। ऑक्सीजन ($2p^4$, प्रतिलोमित, \termsym{3}{P}{2}
सबसे नीचे, फिर $J = 1$, \qty{158.3}{cm^{-1}} पर, और $J = 0$, \qty{227.0}{cm^{-1}} पर):
$158.3/68.7 = 2.30$। नियम 20\,\% के भीतर लागू होता है: \omterm{def:b3:many-electron-atoms:russell-saunders}{रसेल--सॉन्डर्स युग्मन}
हल्के परमाणुओं का अच्छा वर्णन है, यथार्थ नहीं।
\end{example}

\begin{omfigure}
\begin{tikzpicture}[font=\small]
  \draw[omstate] (0,0) -- (1.6,0) node[midway, below] {$3s\ {}^2$S$_{1/2}$};
  \draw[omstate] (0,3.0) -- (1.6,3.0);
  \node[anchor=east] at (-0.1,3.0) {$3p$};
  \draw[omstate] (2.8,2.75) -- (4.4,2.75) node[right] {${}^2$P$_{1/2}$};
  \draw[omstate] (2.8,3.25) -- (4.4,3.25) node[right] {${}^2$P$_{3/2}$};
  \draw[densely dotted, gray] (1.6,3.0) -- (2.8,2.75) (1.6,3.0) -- (2.8,3.25);
  \draw[omfluo] (3.3,2.73) -- (3.3,0.02) node[pos=0.55, left] {D$_1$};
  \draw[omfluo] (3.9,3.23) -- (3.9,0.02) node[pos=0.55, right] {D$_2$};
  \draw[omstate] (2.8,0) -- (4.4,0);
  \draw[densely dotted, gray] (1.6,0) -- (2.8,0);
  \draw[decorate, decoration={brace, amplitude=3pt}] (5.5,3.27) -- (5.5,2.73)
     node[midway, right=3pt, align=left] {\footnotesize \qty{17.2}{cm^{-1}}\\[-1pt]\footnotesize (पैमाने पर नहीं)};
\end{tikzpicture}

{\small सोडियम की D रेखाएँ। $3p$ स्तर \omterm{def:b3:many-electron-atoms:spin-orbit}{चक्रण--कक्षक युग्मन} द्वारा
${}^2$P$_{1/2}$ और ${}^2$P$_{3/2}$ में विपाटित होता है; मूल $3s$ स्तर पर उत्सर्जन से
दो पीली रेखाएँ मिलती हैं, D$_1$, \qty{589.76}{nm} पर, और D$_2$, \qty{589.16}{nm} पर (निर्वात
तरंगदैर्घ्य)।}
\end{omfigure}

\begin{example}[सोडियम द्विक]\label{ex:b3:many-electron-atoms:sodium}
% ledger: lev:Na.3p1_2, lev:Na.3p3_2
सोडियम के $3p$ स्तर 16956.17 और \qty{16973.37}{cm^{-1}} पर हैं: D
रेखाएँ निर्वात में $10^7/16956.17 = \qty{589.76}{nm}$ और \qty{589.16}{nm} पर हैं,
\qty{17.20}{cm^{-1}} के अंतर पर। ${}^2$P के लिए $E(\frac32) - E(\frac12) =
\frac32A$, अतः $A = \qty{11.47}{cm^{-1}}$।
\end{example}

\begin{proposition}[परमाणुओं के लिए वरण नियम]\label{prop:b3:many-electron-atoms:selection-atoms}
\omterm{def:b3:many-electron-atoms:russell-saunders}{रसेल--सॉन्डर्स युग्मन} में विद्युत-द्विध्रुव संक्रमण के लिए आवश्यक है
$\Delta S = 0$, $\Delta L = 0, \pm1$, $\Delta J = 0, \pm1$ (परंतु $J = 0 \to J
= 0$ नहीं), और समता का परिवर्तन: एक-इलेक्ट्रॉन छलाँग के लिए $\Delta l = \pm1$।
\end{proposition}
\begin{proof}[स्थिति]
$\Delta S = 0$ को \cref{ch:b3:electronic-spectroscopy} में सिद्ध किया गया है: द्विध्रुव
\omterm{def:b3:quantum-model-systems:operator}{संकारक} चक्रण पर क्रिया नहीं करता। बाकी द्विध्रुव \omterm{def:b3:quantum-model-systems:operator}{संकारक} के सदिश स्वरूप से
निकलते हैं और यहाँ स्वीकार किए गए हैं।
\end{proof}

\begin{omfigure}
\begin{tikzpicture}
\begin{axis}[width=0.62\linewidth, height=6.2cm, xmin=-0.5, xmax=3.5,
  ymin=158, ymax=173, axis y line=left, axis x line=none,
  ylabel={ऊर्जा / $10^3$ cm$^{-1}$}, unbounded coords=jump, clip=false,
  tick label style={font=\small}, label style={font=\small}]
\addplot[omstate] table[x=x, y=E] {figdata/out/bachelor-3-atomic-levels-helium.dat};
\node[font=\small, anchor=west] at (axis cs:1.05,166.277) {$1s2s\ {}^1$S};
\node[font=\small, anchor=west] at (axis cs:1.05,171.135) {$1s2p\ {}^1$P};
\node[font=\small, anchor=west] at (axis cs:3.05,159.856) {$1s2s\ {}^3$S};
\node[font=\small, anchor=west] at (axis cs:3.05,169.087) {$1s2p\ {}^3$P};
\node[font=\small, anchor=south] at (axis cs:0.5,172.3) {एकल};
\node[font=\small, anchor=south] at (axis cs:2.5,172.3) {त्रिक};
\draw[omfluo] (axis cs:0.5,171.0) -- (axis cs:0.5,166.42) node[pos=0.5, left, font=\scriptsize] {\qty{2059}{nm}};
\draw[omfluo] (axis cs:2.5,168.95) -- (axis cs:2.5,160.0) node[pos=0.5, left, font=\scriptsize] {\qty{1083}{nm}};
\draw[densely dotted, gray] (axis cs:-0.3,159.856) -- (axis cs:2,159.856);
\draw[<->, omThm] (axis cs:-0.2,159.856) -- (axis cs:-0.2,166.277) node[pos=0.5, right, font=\scriptsize] {$2K$};
\end{axis}
\end{tikzpicture}

{\small दो ``हीलियम''। विन्यासों
$1s2s$ और $1s2p$ के एकल और त्रिक स्तर, मूल अवस्था से ऊपर उनकी मापी गई ऊर्जाओं पर
(${}^3$P पद अपने स्तरों के माध्य पर)। रेखाएँ केवल एक ही कुल की
अवस्थाओं को जोड़ती हैं; हर त्रिक अपने एकल से $2K$ नीचे है।}
\end{omfigure}

\begin{history}[दो हीलियम और अपवर्जन सिद्धांत]
हीलियम का नाम 1868 में सूर्य के नाम पर रखा गया और विलियम रैम्ज़े ने 1895 में
उसे पृथ्वी पर पाया। उसकी रेखाओं की दो श्रेणियों ने स्पेक्ट्रमिकीविदों को तीस
वर्षों तक उलझाए रखा। 1925 में वोल्फ़गांग पाउली ने प्रस्तावित किया कि कोई भी दो इलेक्ट्रॉन
समान क्वांटम संख्याएँ साझा नहीं करते, उसी वर्ष जिसमें जॉर्ज
उलेनबेक और सैमुएल गाउडस्मिट ने इलेक्ट्रॉन चक्रण प्रस्तावित किया; 1926 में वर्नर हाइज़ेनबर्ग ने दिखाया कि
केवल तरंग फलन की सममिति ही हीलियम को उसके एकल और
त्रिक कुलों में विभाजित कर देती है।
\end{history}

\begin{inthelab}[सोडियम द्विक का विभेदन]
एक सोडियम लैंप को ग्रेटिंग स्पेक्ट्रममापी की झिरी के सामने रखा जाता है। 600 रेखाएँ
प्रति मिलीमीटर वाली ग्रेटिंग के साथ, दोनों रेखाओं के प्रथम-कोटि कोण
लगभग \num{0.02}\textdegree{} से भिन्न होते हैं: अच्छा स्पेक्ट्रममापी उन्हें अलग कर देता है,
विद्यार्थी प्रिज़्म नहीं। लैंप गर्म चलता है; उसे ठंडा होने के बाद ही
छुआ जाता है।
\end{inthelab}

\section{अभ्यास}

\begin{exercise}[$\star$]\label{exo:b3:many-electron-atoms:1}
लीथियम की मूल अवस्था, $1s^22s$, का \omterm{def:b3:many-electron-atoms:slater}{स्लेटर सारणिक} लिखिए, जिसमें
$2s$ इलेक्ट्रॉन का चक्रण $\alpha$ हो। $1s^3$ के लिए कोई सारणिक क्यों
नहीं होता?
\end{exercise}

\begin{exercise}[$\star$]\label{exo:b3:many-electron-atoms:2}
विन्यासों $p^2$, $p^3$ और $d^2$ के सारणिक गिनिए। $p^3$ के पद
${}^4$S, ${}^2$D, ${}^2$P हैं और $d^2$ के ${}^3$F, ${}^3$P, ${}^1$G,
${}^1$D, ${}^1$S: दोनों गिनतियाँ जाँचिए।
\end{exercise}

\begin{exercise}[$\star$]\label{exo:b3:many-electron-atoms:3}
N, O, F, \ce{Ti^2+} और \ce{Fe^3+} ($d^5$) के मूल पद और स्तर दीजिए।
\end{exercise}

\begin{exercise}[$\star$]\label{exo:b3:many-electron-atoms:4}
पदों ${}^2$D और ${}^3$P के स्तरों को उनकी \omterm{def:b3:quantum-model-systems:degenerate}{अपभ्रष्टताओं} $2J+1$ सहित सूचीबद्ध कीजिए,
और जाँचिए कि \omterm{def:b3:quantum-model-systems:degenerate}{अपभ्रष्टताओं} का योग $(2L+1)(2S+1)$ है।
\end{exercise}

\begin{exercise}[$\star\star$]\label{exo:b3:many-electron-atoms:5}
कार्बन के ${}^3$P स्तर 0, 16.4 और \qty{43.4}{cm^{-1}} पर हैं। अंतरालों का
अनुपात और हर अंतराल से चक्रण--कक्षक स्थिरांक $A$ परिकलित कीजिए।
$A$ के दोनों मानों का अंतर क्या बताता है?
\end{exercise}

\begin{exercise}[$\star\star$]\label{exo:b3:many-electron-atoms:6}
सोडियम के $3p$ स्तर 16956.17 और \qty{16973.37}{cm^{-1}} पर हैं।
D रेखाओं के निर्वात तरंगदैर्घ्य, \unit{cm^{-1}} और
\unit{meV} में उनका अंतर, और $A$ परिकलित कीजिए, $3p$ पद के लिए।
\end{exercise}

\begin{exercise}[$\star\star$]\label{exo:b3:many-electron-atoms:7}
सोडियम के इनमें से कौन-से संक्रमण उत्सर्जन में अनुमत हैं: $3p \to 3s$,
$3d \to 3s$, $3d \to 3p$, $4s \to 3s$, $4s \to 3p$? हर एक का औचित्य दीजिए।
\end{exercise}

\begin{exercise}[$\star\star$]\label{exo:b3:many-electron-atoms:8}
% ledger: lev:He.2p3P2, lev:He.2p3P1, lev:He.2p3P0, lev:He.2p1P1
हीलियम के $1s2p$ स्तर \termsym{3}{P}{2}, \termsym{3}{P}{1},
\termsym{3}{P}{0} हैं, 169086.76, 169086.84, \qty{169087.83}{cm^{-1}} पर, और
\termsym{1}{P}{1}, \qty{171134.89}{cm^{-1}} पर। ${}^3$P पद की माध्य ऊर्जा
परिकलित कीजिए, फिर $K(1s,2p)$ eV में। $K(1s,2s) = \qty{0.398}{eV}$ से तुलना कीजिए
और अंतर समझाइए।
\end{exercise}

\begin{exercise}[$\star\star$]\label{exo:b3:many-electron-atoms:9}
इस अध्याय की विधि से $d^2$ के पद ज्ञात कीजिए, सबसे बड़े $M_L$ से
आरंभ करते हुए (पूरी सारणी लिखना छोड़ सकते हैं, परंतु हर पद का औचित्य दीजिए)।
\end{exercise}

\begin{exercise}[$\star\star\star$]\label{exo:b3:many-electron-atoms:10}
दिखाइए कि त्रिक आकाशी फलन $\frac{1}{\sqrt2}[1s(1)2s(2) -
2s(1)1s(2)]$, $\mathbf r_1 = \mathbf r_2$ होने पर लुप्त होता है और एकल
फलन नहीं। इसे $E_+ - E_-$ के चिह्न से जोड़िए।
\end{exercise}

\begin{exercise}[$\star\star\star$]\label{exo:b3:many-electron-atoms:11}
% ledger: lev:Ni2.3F3, lev:Ni2.3F2
\ce{Ni^2+} ($d^8$) के लिए ${}^3$F स्तर 0 ($J = 4$), 1360.7 ($J = 3$)
और \qty{2269.6}{cm^{-1}} ($J = 2$) पर हैं। क्रम की व्याख्या कीजिए, अंतरालों का
अनुपात परिकलित कीजिए और लांडे की भविष्यवाणी से तुलना कीजिए। बड़े अंतराल से
$A$ दीजिए।
\end{exercise}

\begin{exercise}[$\star\star\star$]\label{exo:b3:many-electron-atoms:12}
दिखाइए कि $\hat{\mathbf L}\cdot\hat{\mathbf S}$ का \omterm{def:b3:quantum-model-systems:hermitian}{प्रत्याशा मान}, किसी पद की
सभी $(2L+1)(2S+1)$ अवस्थाओं पर जोड़ने पर, शून्य है, अतः \omterm{def:b3:many-electron-atoms:spin-orbit}{चक्रण--कक्षक
युग्मन} स्तरों के भारित माध्य को अपरिवर्तित छोड़ता है। इसे
कार्बन के ${}^3$P पर जाँचिए।
\end{exercise}

\section{समस्या: दो हीलियम}

\begin{problem}[{सप्ताहांत समस्या --- हीलियम के एकल और त्रिक कुल:
सारणिक, संक्रमण-धातु आयनों के पद, दोनों कुलों को अलग रखने वाले वरण
नियम, और मापा गया विनिमय समाकल}]
\label{pb:b3:many-electron-atoms:1}
% ledger: lev:He.2s3S1, lev:He.2s1S0, lev:He.2p3P2, lev:He.2p3P1, lev:He.2p3P0, lev:He.2p1P1, lev:Ti2.3F3, lev:Ti2.3F4, lev:Ti2.3P0, lev:Ti2.3P1, lev:Ti2.3P2
आँकड़े (NIST, \unit{cm^{-1}} में, हर स्पीशीज़ की मूल अवस्था से ऊपर)। हीलियम:
$1s2s$ \termsym{3}{S}{1} 159855.97, \termsym{1}{S}{0} 166277.44; $1s2p$
\termsym{3}{P}{2,1,0} 169086.76, 169086.84, 169087.83, \termsym{1}{P}{1}
171134.89। \ce{Ti^2+}: \termsym{3}{F}{2,3,4} 0, 184.9, 420.4; \termsym{3}{P}{0,1,2}
10538.4, 10603.6, 10721.2। $\qty{1}{eV} = \qty{8065.54}{cm^{-1}}$।

\medskip\noindent\textbf{भाग I --- \omterm{def:b3:many-electron-atoms:spin-orbital}{चक्रण-कक्षक} और सारणिक।}
\begin{enumerate}
  \item विन्यास $1s2s$ के लिए उपलब्ध चार \omterm{def:b3:many-electron-atoms:spin-orbital}{चक्रण-कक्षक} सूचीबद्ध कीजिए।
  \item वे चार सारणिक लिखिए जिनमें एक इलेक्ट्रॉन $1s$ में और एक
        $2s$ में हो।
  \item दिखाइए कि $|1s\alpha\,2s\alpha|$ एक प्रतिसममित
        आकाशी फलन और चक्रण फलन $\alpha(1)\alpha(2)$ का गुणनफल है।
  \item $|1s\alpha\,2s\beta|$ और $|1s\beta\,2s\alpha|$ को संयोजित करके उसी आकाशी फलन का
        $M_S = 0$ घटक, और एक चौथी अवस्था, प्राप्त कीजिए।
  \item कौन-सा आकाशी फलन, सममित या प्रतिसममित,
        $S = 0$ के साथ जाता है? $S = 1$ के साथ कौन-सा?
  \item इन्हें एकल और त्रिक क्यों कहा जाता है?
\end{enumerate}

\medskip\noindent\textbf{भाग II --- \ce{Ti^2+} के पद।}
\begin{enumerate}[resume]
  \item \ce{Ti^2+} का विन्यास और उसके
        सारणिकों की संख्या दीजिए।
  \item हुंड के नियमों से उसका मूल पद और स्तर ज्ञात कीजिए; आँकड़ों से
        जाँचिए।
  \item उसके अन्य पद ${}^3$P, ${}^1$G, ${}^1$D और ${}^1$S हैं: अवस्थाओं की गिनती
        जाँचिए।
  \item क्या ${}^3$F की \omterm{def:b3:many-electron-atoms:spin-orbit}{सूक्ष्म संरचना} सामान्य है या प्रतिलोमित? लांडे के
        अंतराल नियम का परीक्षण कीजिए।
  \item ${}^3$F और ${}^3$P पदों की भारित माध्य ऊर्जाएँ परिकलित कीजिए।
  \item उनका अंतर $15B$ है, जहाँ $B$ आयन का एक राका
        प्राचल है (\cref{ch:b3:complex-spectra-magnetism})। $B$ परिकलित कीजिए।
\end{enumerate}

\medskip\noindent\textbf{भाग III --- दो कुल क्यों।}
\begin{enumerate}[resume]
  \item कौन-सा वरण नियम किसी एकल और किसी
        त्रिक के बीच संक्रमण को वर्जित करता है?
  \item निम्नतम त्रिक, $1s2s$ \termsym{3}{S}{1}, मूल
        अवस्था $1s^2$ \termsym{1}{S}{0} पर उत्सर्जन नहीं कर सकता। दो कारण दीजिए। ऐसी अवस्था
        को क्या कहते हैं?
  \item रेखा $1s2p\ {}^3$P $\to 1s2s\ {}^3$S का तरंगदैर्घ्य परिकलित कीजिए
        (${}^3$P स्तरों का माध्य लीजिए)।
  \item $1s2p\ {}^1$P $\to 1s2s\ {}^1$S का तरंगदैर्घ्य परिकलित कीजिए।
  \item $1s2p\ {}^1$P $\to 1s^2\ {}^1$S का तरंगदैर्घ्य परिकलित कीजिए। यह स्पेक्ट्रम के
        किस क्षेत्र में है?
  \item समझाइए कि उन्नीसवीं शताब्दी के स्पेक्ट्रमिकीविद् ये रेखाएँ देखकर
        दो भिन्न गैसों में क्यों विश्वास करते थे।
\end{enumerate}

\medskip\noindent\textbf{भाग IV --- मापा गया \omterm{def:b3:many-electron-atoms:exchange}{विनिमय समाकल}।}
\begin{enumerate}[resume]
  \item $1s2s$ एकल और त्रिक की प्रथम-कोटि ऊर्जाएँ
        $E_{1s}$, $E_{2s}$, $J$ और $K$ के पदों में लिखिए।
  \item कौन-सी नीचे है, और क्यों, इलेक्ट्रॉनों की स्थितियों के संदर्भ में?
  \item $1s2s$ एकल--त्रिक अंतराल \unit{cm^{-1}} और eV में परिकलित कीजिए।
  \item $K(1s,2p)$ को eV में परिकलित कीजिए, $1s2p$ स्तरों से।
  \item $K(1s,2s)$, $K(1s,2p)$ से बड़ा क्यों है?
  \item क्या हीलियम की उत्तेजित अवस्थाएँ हुंड के पहले नियम की पुष्टि करती हैं?
  \item परिणाम बताइए: हीलियम का \omterm{def:b3:many-electron-atoms:exchange}{विनिमय समाकल} $K(1s,2s)$,
        eV में।
\end{enumerate}
\end{problem}
